\chapter{Foundations for Verification}
\label{ch:verification}

\section{Emitter correctness}

For every admitted function shape, a correctness theorem states that
running the emitted Lean definition agrees with evaluating the
\texttt{CoreIR} function: semantics in, same semantics out, on both
the ok and the error paths (each shape ships ok- and error-path
corollaries alongside the main agreement). These proofs
go by structural or fuel induction, and each new operation earns its
per-operation lemmas before the gate admits it. Closed entries take a
shorter path: both sides are closed terms, so kernel-checked evaluation
decides them, replacing hand-written evaluation proofs. The result is that trust
in any emitted program reduces to trust in the evaluator plus the
per-shape proof --- never to the output text itself, which is additionally
pinned byte-for-byte by golden files.

\section{Validator and disciplines}

The validator is the project's bouncer: it maps admitted corpus shapes to
their canonical functions and rejects everything else with a precise,
actionable code (aliasing, escape, possible out-of-bounds, out-of-subset).
Around it sit the small disciplines that keep the common case pure:
verified noalias derivation from CIR text, a read-only sharing discipline
for recovered readers, and extracted scope/alloca lifetime evidence that
informs admission messages, pins residue reports, and feeds admission
pins. Each discipline follows the same rule as everything else: admit
exactly what is understood, reject the rest loudly.

\section{Specifications and tactics}

Users prove properties as pure equations over emitted code, in the style
of the gallery checked into \texttt{Circe.Specs}: sortedness of a fill
loop, first-match minimality of a search, prefix preservation of a
reallocation. The sorting case studies carry the full contracts: each
admitted sort returns a sorted permutation of its input. Proof
ergonomics come from a small tactic family. One
simplifier set bundles checked-operation unfoldings, range bridges, and
whole-program equations; fuel automation discharges the arithmetic that
loop proofs otherwise repeat; a selector tactic closes the borrow-return
lens laws; a registered step cascade discharges program-evaluation
steps, with composer fuel splits packed into a single obtain; and a
fold-end tactic closes the fold-termination script shared across all
bounded folds. The emitter also scaffolds unverified spec stubs next to each
generated file --- signature, body reference, edge cases, and a
property-test entry --- so new programs arrive with their proof
obligations already sketched. The standing audit result is that everyday
proofs are simplifier-first with at most a couple of manual steps.

\section{Hoare-style verification}

A second verification style grows beside the equational one: Hoare
triples over the memory evaluator. A triple $\{P\}~s~\{Q\}$ here is
thin sugar, not new semantics. \texttt{HTriple} states that statement
\texttt{s}, run from any state satisfying the pre-predicate \texttt{P},
evaluates --- at any fuel --- to an \texttt{.ok} outcome satisfying the
post-predicate \texttt{Q}. States are triples of environment, memory,
and layout; outcomes distinguish falling through from returning a
value, among others. Correctness is total over \texttt{.ok} only:
failure is excluded by construction, and every way a program could fail
(overflow, a failed store, a missed guard) is discharged as an explicit
hypothesis instead --- debt in the open, matching the transfer style,
never an axiom. A program-level twin covers loop bodies with calls,
which the call-free triple cannot see.

Fuel deserves one honest paragraph, because the loop rule turns on it.
The evaluator is fuel-bounded, so loop invariants are indexed by fuel:
\texttt{I k} describes the states that may arise with \texttt{k} steps
left, and fuel itself does the bounding, which means no separate
termination measure is needed for soundness. The index is load-bearing
rather than decorative: at fuel zero a true-guard loop errors, so
\texttt{I 0} admits only guard-false states, and each true-guard step
must deliver a state satisfying \texttt{I k} from one satisfying
\texttt{I (k+1)}. Bodies that break or return out of the loop are
excluded for now; the handler arms exist, but the exit-post plumbing is
deferred with the gap named, not silently assumed.

The rules themselves are direct wrappers over the existing memory
evaluator equations, one per statement shape: skip, sequencing with the
mid-outcome pinned to falling through, consequence for strengthening
pre-conditions and weakening post-conditions, both conditional branches
over plain boolean guards, return with its value as the outcome, and
let/assign for environment extension and update. At program level, call
rules reduce the caller to the callee run, carrying caller memory
across definitionally. Storage gets the
interesting rule: a store produces the stored-cell fact and preserves
one framed cell across the write --- the untouched cell carried through
as an explicit conjunct. The store index is pinned to zero, which is
complete rather than limiting, because the cell assertion describes a
singleton one-word block and anything else would be out of bounds.
Loads appear only inside store value expressions, through an
expression-level helper, since loads have no value-side mirror by
design.

Confidence comes from a gallery of closed-instance proofs that double
as regression tests: a two-cell increment with framing, both arms of a
conditional writer over its real function body, a countdown loop, a
pair-returning conditional, a read-only accumulation loop with
per-prefix arithmetic debt, and an in-place increment loop carrying
the evolving prefix in its invariant. Each is proved end to end from
the rules with its debts as hypotheses, each ships with a runtime
check, and the driver runs them all. Whole-list assertions generalize
the singleton cell facts wherever loops range over blocks. Sequencing
past an early return remains the named open lemma on the memory side;
until it lands, straight-line bodies return last, and the chapter will
say so for exactly as long as it stays true.
